\section{From high traces to the graph estimate}\label{sec:localization}

We now prove \Cref{prop:graph-estimate} from \Cref{thm:high-trace}.
The latter controls averages over a finite residue space. The present
step converts those averages into a bound for arbitrary test sequences
on long intervals, including complex coefficients $a_d$.

Choose a power of two $D_0$ with
\[
 B^{10}H\le D_0<2B^{10}H,
\]
and partition the positive integers into blocks
$I_k=\{kD_0+1,\ldots,(k+1)D_0\}$, $k\ge0$.
For large $B$, all primes in $\mathcal S$ are odd, so
$\gcd(D_0,Q_B)=1$. Let $Y$ be the periodic vertex set from
\Cref{thm:high-trace}. On the coordinates of $I_k$, define the matrix
\[
 M_k(x,y)=
 \begin{cases}
 \displaystyle\1_Y(x)\1_Y(y)
 \frac{a_dg_d(x,y)}{\sqrt{W(x)W(y)}},
       &y=x+hd,\quad d\in\mathcal D,\\[5pt]
 0,&\text{otherwise}.
 \end{cases}
\]
This is a directed matrix and need not be self-adjoint. We therefore
use a moment of its singular values.

\subsection{Averaging the singular-value moment}

Set $m=\ell/2$. With $y_m=y_0$, matrix multiplication gives exactly
\begin{equation}\label{eq:localization-trace}
 \operatorname{Tr}\bigl((M_k^*M_k)^m\bigr)
 =\sum_{y_0,\ldots,y_{m-1}\in I_k}\ 
   \sum_{x_1,\ldots,x_m\in I_k}
   \prod_{j=1}^m\overline{M_k(x_j,y_{j-1})}M_k(x_j,y_j).
\end{equation}
A nonzero term describes the closed line
\[
 y_0,x_1,y_1,x_2,y_2,\ldots,x_m,y_0,
\]
with alternating negative and positive steps of lengths $hd_i$.
Its scalar coefficient is
$\prod_{j=1}^m\overline{a_{d_{2j-1}}}a_{d_{2j}}$, of modulus at most one.
Symmetry and reality of $g_d$ identify its edge numerator with that of
$K_{\mathbf d}$ in \eqref{eq:trace-kernel}. For the successive vertices
$v_0,\ldots,v_\ell=v_0$, its denominator satisfies
\[
 \prod_{i=1}^{\ell}\sqrt{W(v_{i-1})W(v_i)}
       =\prod_{i=0}^{\ell-1}W(v_i).
\]
Repeated vertices are counted with multiplicity in this identity.
The indicator factors reduce to $\prod_{i=0}^\ell\1_Y(v_i)$.

Fix the initial position $r\in\{1,\ldots,D_0\}$ and write
$v_i=kD_0+r+b_i$. The block restrictions are precisely
$1\le r+b_i\le D_0$ for all $i$; they depend only on $r$ and the
ordered line. As $k$ runs modulo $Q_B$, the starting point $kD_0+r$
is uniform modulo $Q_B$. Averaging \eqref{eq:localization-trace} over
this period, and taking absolute values after each line's residue
average, gives
\begin{align}
 0\le\frac1{Q_B}\sum_{k\bmod Q_B}
 \operatorname{Tr}\bigl((M_k^*M_k)^m\bigr)
 &\le D_0\sum_{\mathbf d}
       \left|\E K_{\mathbf d}(n)
             \prod_{i=0}^{\ell}\1_Y(n+b_i)\right|\notag\\
 &\le D_0 B^{-(1+c_*)\ell}.\label{eq:localization-moment}
\end{align}
A starting position and its ordered signed line determine all indices
in \eqref{eq:localization-trace}, so no further multiplicity occurs.

Call a block exceptional if
$\|M_k\|>B^{-1-c_*/2}$, where the norm is the Euclidean operator norm.
Since $\|M_k\|^\ell\le\operatorname{Tr}((M_k^*M_k)^m)$, its fraction
$\delta$ among one period of blocks satisfies
\begin{equation}\label{eq:exceptional-blocks}
 \delta\le D_0B^{-c_*\ell/2}
       =\exp\bigl(-\Omega(B\log B)\bigr).
\end{equation}
Here $\log D_0\le C_0B+10\log B+O(1)$ and $\ell\sim2B$.
The threshold for this estimate may depend on $C_0$.

\subsection{Discarded edges}

We must control the weight on deleted vertices, exceptional blocks and
block boundaries before applying the operator norm. A common absolute
degree bound handles all three. Dropping cutoffs and extending the
sum to every squarefree product of primes in $\mathcal S$ gives,
separately for either sign,
\begin{align}
 \sum_{d\in\mathcal D}|g_d(x,x\pm hd)|
 &\le\prod_{p\in\mathcal C}(1+A\1_{p\mid x})
    \prod_{p\in\mathcal Z}
       \left(1+\left|\1_{p\mid x}-\frac\theta p\right|\right)
 \le W(x)e^{\theta v_Z}.\label{eq:absolute-degree}
\end{align}
For $p\in\mathcal Z$ dividing $x$, its factor is $2-\theta/p\le2$;
otherwise it is at most $e^{\theta/p}$. The same bound applies to
incoming edges at their terminal vertex, since the residues of both
endpoints agree at every prime in their divisor label.

All parameters are now fixed. Periodic averages over the integers equal
the corresponding residue averages as $X\to\infty$. By
\Cref{thm:high-trace} and \eqref{eq:weight-moments}, the total absolute
edge weight incident to $Y^c$, divided by $X$, has limsup at most
\[
 2e^{\theta v_Z}\E\bigl(W\1_{Y^c}\bigr)
 \le 2e^{\theta v_Z}(\E W^2)^{1/2}\Pbb(Y^c)^{1/2}
 \le B^{O_A(1)}P_0^{-1/2+o(1)}.
\]
The fixed extension of the interval to $X+\tau hH$ does not affect this
limsup. The last expression is smaller than every fixed negative power
of $B$.

Let $E$ be the union of exceptional blocks. It has period $D_0Q_B$ and
density $\delta$. Cauchy--Schwarz over that full period gives
\[
 \lim_{X\to\infty}\frac1X\sum_{x\le X}W(x)\1_E(x)
 \le (\E W^2)^{1/2}\delta^{1/2}.
\]
There is no independence assertion between the exceptional-block event
and the prime residues. By \eqref{eq:absolute-degree}, the within-block
terms in $E$ cost at most
\[
 e^{\theta v_Z}(\E W^2)^{1/2}\delta^{1/2}
 \le B^{O_A(1)}D_0^{1/2}B^{-c_*\ell/4}
 =\exp\bigl(-\Omega(B\log B)\bigr)
\]
after normalization by $X$ and passage to the limsup.

A forward edge crossing a block boundary begins in the last
$\lceil\tau hH\rceil$ positions of its block. Position modulo $D_0$
and residue modulo $Q_B$ are independent over a full period $D_0Q_B$.
Consequently the crossing cost is at most
\[
 \frac{\lceil\tau hH\rceil}{D_0}L_0e^{\theta v_Z}
 \ll_h B^{-10}L_0e^{\theta v_Z}
 =o\bigl(L_0B^{-1-c_*/2}\bigr).
\]
In the first two estimates we may also compare with this target, since
$L_0\ge1$. Thus all three discarded contributions are negligible.

\subsection{The bilinear estimate}

Use the standard complex inner product, conjugate-linear in its first
argument, and define
\[
 u(x)=\overline{F(x)}\sqrt{W(x)}\1_{1\le x\le X},\qquad
 v(y)=G(y)\sqrt{W(y)}\1_{1\le y\le X+\tau hH}.
\]
Then $\langle u|_{I_k},M_kv|_{I_k}\rangle$ is exactly the desired
unnormalized sum over edges in $I_k$ with endpoints in $Y$.
On nonexceptional blocks, Cauchy--Schwarz in the block index yields
\begin{align*}
 \left|\sum_{k\ \mathrm{nonexceptional}}
      \langle u|_{I_k},M_kv|_{I_k}\rangle\right|
 &\le B^{-1-c_*/2}
       \left(\sum_{x\le X}W(x)\right)^{1/2}
       \left(\sum_{y\le X+\tau hH}W(y)\right)^{1/2}\\
 &=(L_0+o_{X\to\infty}(1))X B^{-1-c_*/2}.
\end{align*}
The source cutoff in $u$ also handles the last partially used block.
Adding back the three negligible edge classes proves
\Cref{prop:graph-estimate}. This closes the deferred graph input in
\Cref{sec:analytic-transfer}, and hence completes the proof of
\Cref{thm:elliott}.
